% Chapter 6 - Testing, Evaluation and Conclusion (rebuilt for the P2 Mid report)
\chapter{Testing, Evaluation and Conclusion}

This chapter describes how Pakalorie is tested and evaluated against the requirements in Chapter 3 and the
objectives in Chapter 1. It first explains the testing strategy and the test cases, with the result of each test
that has been run. It then reports how model quality, speed and usability are measured, and ends with the
limitations of the system, the conclusion and future work. Where a test has not been run, the table says so.

\section{Testing Strategy}
Testing follows the usual levels of software testing~\cite{test_plan_1,system_test}. Because the project uses the
Agile Incremental model, tests are run at the end of every milestone and repeated on earlier modules to catch
broken features early~\cite{unit_testing}.
\begin{itemize}[noitemsep]
  \item \textbf{Unit testing:} single functions are tested alone, for example the calorie compensation rule, the
        burn estimate and the food service. The backend has a unit test suite for the calorie engine, the depth
        service, the food service and the seed data.
  \item \textbf{Integration testing:} modules are tested together, for example app to API to database. The backend
        has a live-database integration suite.
  \item \textbf{System testing:} the full flow is tested on a device, from sign-up to a saved meal.
        Black-box test cases are written from the functional requirements~\cite{integrate_test}.
  \item \textbf{Performance testing:} response time of the recognition endpoint.
  \item \textbf{Usability testing:} planned, not yet carried out.
  \item \textbf{Model evaluation:} the model is tested on images it has not seen during training.
\end{itemize}

\section{Test Environment}
\begin{itemize}[noitemsep]
  \item \textbf{Mobile devices:} an iPhone 17 Pro simulator running Expo Go 54. A real Android phone and a real
        iPhone have not yet been used, so camera tests are pending.
  \item \textbf{Backend:} the Docker setup described in Chapter 5, reached over its public HTTPS address.
  \item \textbf{Data:} gallery photos of dishes, test accounts, and a validation set of food images.
  \item \textbf{Languages:} English only at present.
\end{itemize}

\section{Functional Test Cases}
Table~\ref{tab:tc} lists the main test cases. Each case is linked to a functional requirement (FR) from
Table~3.1. The Status column uses these words: \emph{Pass}, \emph{Partial} (part of the case was checked),
\emph{Not run}, \emph{Not built} (the feature does not exist yet) and \emph{Not met}. The results come from one test
session on 11 November 2026 on the simulator unless a note says otherwise.

\begin{longtable}{|l|l|>{\raggedright\arraybackslash}p{3.7cm}|>{\raggedright\arraybackslash}p{3.2cm}|>{\raggedright\arraybackslash}p{2.6cm}|}
\caption{Functional Test Cases}\label{tab:tc}\\
\hline
\textbf{ID} & \textbf{FR} & \textbf{Test Case} & \textbf{Expected Result} & \textbf{Status} \\ \hline
\endfirsthead
\multicolumn{5}{c}{\tablename\ \thetable{} -- Continued from previous page}\\ \hline
\textbf{ID} & \textbf{FR} & \textbf{Test Case} & \textbf{Expected Result} & \textbf{Status} \\ \hline
\endhead
\hline \multicolumn{5}{r}{Continued on next page}\\
\endfoot
\hline
\endlastfoot
TC01 & FR1 & Register with a valid email and password & Account is created and a profile document is stored & Pass \\ \hline
TC02 & FR1 & Register with an invalid email or weak password & Clear error message; no account is created & Not run \\ \hline
TC03 & FR2 & Log in with correct email and password & Home screen opens & Not run \\ \hline
TC04 & FR2 & Sign in with Google & Home screen opens & Not run (needs a development build) \\ \hline
TC05 & FR3 & Change weight and calorie goal & New values are saved and shown after restart & Not run \\ \hline
TC06 & FR4 & Open the app after 24 hours of inactivity & Login screen is shown & Not run \\ \hline
TC07 & FR5 & Save a meal after recognition & Meal appears under the user's ID in Firestore & Pass \\ \hline
TC08 & FR6 & Capture a food photo with the camera & Image is sent and a result is shown & Not run (simulator has no camera) \\ \hline
TC09 & FR7 & Choose a food photo from the gallery & Image is sent and a result is shown & Pass (haleem, 90--93\%) \\ \hline
TC10 & FR8 & Submit a photo on 4G or WiFi & Result screen appears within 3 seconds & Not met on the Claude path (about 6 s); met on the YOLOv8 path (0.8--2.2 s) \\ \hline
TC11 & FR9 & Open the result screen & Calories, protein, fat and carbohydrates are shown & Pass \\ \hline
TC12 & FR10 & Open History and the weekly trend view & Meals are listed newest first; trends match the saved data & Partial (Macro Split matched; history list not run) \\ \hline
TC13 & FR11 & Photo of a single known dish & Correct label is returned & Pass (one dish tested) \\ \hline
TC14 & FR11 & Photo with several dishes & Each dish is detected as a separate item & Not built (classifier names one dish) \\ \hline
TC15 & FR12 & Open the result of any recognition & A confidence score is shown & Pass \\ \hline
TC16 & FR13 & Photo that gives confidence below 70\% & Alternative food names are suggested & Not run \\ \hline
TC17 & FR14 & Log a known item while offline & Item is logged from cache and synced when online & Not built \\ \hline
TC18 & FR15, FR16 & Call POST /recognize and a food endpoint; call with an unknown identifier & Valid JSON for known items; 404 for unknown & Partial (recognise and calories tested; unknown identifier not run) \\ \hline
TC19 & FR17, FR18 & Photos of the same dish in a small and a large portion & Larger portion gives a higher portion factor and more calories & Not run \\ \hline
TC20 & FR19 & Ask for calories of a dish through the RAG engine & Answer uses the retrieved nutrition values & Pass (database-grounded, 300 kcal for haleem) \\ \hline
TC21 & FR20, FR23 & Daily calories above the goal on earlier days & Adjusted goal is shown with a safe minimum & Partial (goal and activity credit shown; carry-over across days checked by calculation only) \\ \hline
TC22 & FR21 & Run the quantised model on the test set & Labels match the original model closely; speed improves & Pass for size and accuracy (Table~\ref{tab:quant}); not yet on a phone \\ \hline
TC23 & FR22 & Open the dashboard with meals & Charts load and match the saved data & Partial (Macro Split matched 27/34/39\%; empty state shown; other charts not in this run) \\ \hline
TC24 & FR24 & Run live recognition with the camera & Labels update without pressing a button & Partial (toggle and overlay shown; camera not tested; concurrent requests not tested) \\ \hline
TC25 & FR25, FR26 & Switch the language to Urdu & Text and layout become right-to-left; Urdu food names appear & Not built \\ \hline
TC26 & FR27, FR28 & Save a meal with health sync turned on & Nutrition entry appears in Google Fit or HealthKit & Not built \\ \hline
TC27 & FR29 & Open the health dashboard & Nutrition and activity data are shown together & Not built \\ \hline
TC28 & FR30--FR32 & Chat with the coach about the user's day & Personal, simple reply with a disclaimer & Not met (error 400 without an API key) \\ \hline
\end{longtable}

\subsection{Non-Functional Test Cases}
\begin{table}[H]
\centering
\caption{Non-Functional Test Cases}
\label{tab:nf}
\small
\begin{tabularx}{\textwidth}{|l|>{\raggedright\arraybackslash}p{4.3cm}|>{\raggedright\arraybackslash}X|>{\raggedright\arraybackslash}p{2.6cm}|}
\hline
\textbf{ID} & \textbf{Test Case} & \textbf{Expected Result} & \textbf{Status} \\ \hline
NF01 & Call an endpoint with no token or a wrong token & Server answers 401 and does no work & Not built (token check planned) \\ \hline
NF02 & Check the server storage after several recognitions & No food images are stored & Not run \\ \hline
NF03 & Try to reach the API over plain HTTP & Request is refused or redirected to HTTPS & Not run \\ \hline
NF04 & Read another user's documents with a test account & Firestore rules deny access & Not run \\ \hline
NF05 & Send an image larger than the limit & App compresses it before upload; server rejects oversize files & Not run \\ \hline
NF06 & Look at the result screen & Disclaimer is visible & Pass \\ \hline
\end{tabularx}
\end{table}

\section{Model Evaluation}\label{sec:modeleval}
\subsection{Metrics}
The recognition model is a classifier and is evaluated on a validation set that it did not train on. The
following measures are used.
\begin{itemize}[noitemsep]
  \item \textbf{Top-1 accuracy} = correct first predictions / total predictions.
  \item \textbf{Top-5 accuracy} = the share of images where the true class is among the five best guesses.
  \item \textbf{Confusion matrix} shows which classes are mixed up, which guides new data collection.
\end{itemize}
Precision, recall and mAP are measures for detection models with bounding boxes; they do not apply to this
classifier and are not reported.

\subsection{Targets and Results}
\begin{table}[H]
\centering
\caption{Model Evaluation Targets and Measured Values}
\label{tab:modeleval}
\small
\begin{tabularx}{\textwidth}{|>{\raggedright\arraybackslash}p{5.2cm}|>{\raggedright\arraybackslash}p{3.2cm}|>{\raggedright\arraybackslash}X|}
\hline
\textbf{Measure} & \textbf{Target} & \textbf{Measured} \\ \hline
Top-1 accuracy (217 classes, official evaluation) & At least 85\% (Chapter 1) & 58.5\% (top-5 86.6\%); target not met \\ \hline
Multi-item images: items found / items present & Reported & Not applicable (one dish per image) \\ \hline
Accuracy change after INT8 quantisation (static) & Small loss only & $-3.7$ points top-1 and $-0.2$ points top-5 (Table~\ref{tab:quant}) \\ \hline
Model size, original vs.\ static INT8 & Smaller & 6.89~MB vs.\ 1.89~MB (3.6$\times$) \\ \hline
Inference time, original vs.\ static INT8 (CPU, one image) & Within the 3-second goal & 2.4~ms vs.\ 2.1~ms (model only) \\ \hline
\end{tabularx}
\end{table}

\subsection{Portion and Calorie Check}
To check portion estimation, a small set of meals should be weighed on a kitchen scale, photographed, and compared
with the estimate. This check has not been done. The portion estimate is relative and is not expected to give
grams. The calories of one tested dish (haleem, 300~kcal) came from database rows and were not compared with a scale.

\section{Performance Evaluation}
\begin{table}[H]
\centering
\caption{Performance Targets and Measurement Method}
\label{tab:perf}
\small
\begin{tabularx}{\textwidth}{|>{\raggedright\arraybackslash}p{3.3cm}|>{\raggedright\arraybackslash}p{2.6cm}|>{\raggedright\arraybackslash}X|>{\raggedright\arraybackslash}p{3.1cm}|}
\hline
\textbf{Measure} & \textbf{Target} & \textbf{How It Is Measured} & \textbf{Measured} \\ \hline
Time from upload to result & 3 seconds or less & Timed over many runs on 4G and WiFi & YOLOv8 path: about 0.77 s (laptop) and 2.2 s (1~MB photo, another computer). Claude path: about 6 s (one run) \\ \hline
Home screen start-up & 2 seconds or less & Timed from launch on a mid-range phone & Not measured \\ \hline
Concurrent users & At least 50 without slowdown & Load test that simulates 50 users & Not measured \\ \hline
Upload size & 1 MB or less & Size checked after compression & Not measured \\ \hline
Workflow length & 3 interactions or less & Count of taps: open, capture, view result & Met in the tested flow (open Scan, choose a photo, view the result) \\ \hline
Server uptime & At least 99\% & Health check during the evaluation period & Not measured \\ \hline
\end{tabularx}
\end{table}

\section{Usability Evaluation}
Usability is to be tested with a group of users that includes health-conscious people, fitness users and a nutrition
student, like the interview group in Chapter 3. Each person is asked to do five tasks without help: (1) register and
log in, (2) capture a meal and read the calories, (3) save the meal, (4) switch to Urdu, and (5) ask the coach one
question. For each task the team records whether it was completed and how long it took. After the tasks, each person
fills in the System Usability Scale (SUS) questionnaire~\cite{sus}. This study has \emph{not} been carried out yet, so
Table~\ref{tab:usab} reports no results.

\begin{table}[H]
\centering
\caption{Usability Measures}
\label{tab:usab}
\small
\begin{tabularx}{\textwidth}{|>{\raggedright\arraybackslash}p{5cm}|>{\raggedright\arraybackslash}p{4cm}|>{\raggedright\arraybackslash}X|}
\hline
\textbf{Measure} & \textbf{Target} & \textbf{Measured} \\ \hline
Task completion rate & All five tasks & Not yet carried out \\ \hline
Time to capture and read result & Within 3 interactions & Met in the tested flow \\ \hline
SUS score & Reported & Not yet carried out \\ \hline
\end{tabularx}
\end{table}

\section{Requirement Traceability}
Table~\ref{tab:trace} links each module to its requirements, to the test cases that check them, and to the section where
it is built.

\begin{table}[H]
\centering
\caption{Requirement Traceability Matrix}
\label{tab:trace}
\small
\begin{tabularx}{\textwidth}{|>{\raggedright\arraybackslash}p{4.6cm}|l|l|>{\raggedright\arraybackslash}X|}
\hline
\textbf{Module} & \textbf{FRs} & \textbf{Test Cases} & \textbf{Built In} \\ \hline
1. Authentication & FR1--FR5 & TC01--TC07, NF01, NF04 & Sections 5.5.2, 5.6.1 \\ \hline
2. Core Mobile UI & FR6--FR10 & TC08--TC12 & Section 5.6.2 \\ \hline
3. Food Recognition & FR11--FR14 & TC13--TC17 & Sections 5.4, 5.5.4 \\ \hline
4. Food Database API & FR15, FR16 & TC18 & Section 5.5.3 \\ \hline
5. Volume and Depth & FR17, FR18 & TC19 & Section 5.5.5 \\ \hline
6. Calorie Engine and RAG & FR19 & TC20 & Section 5.5.6 \\ \hline
7. Model Optimisation & FR21 & TC22 & Section 5.4.4 \\ \hline
8. Visualisation and Compensation & FR20, FR22, FR23 & TC21, TC23 & Sections 5.5.7, 5.6.3 \\ \hline
9. Real-Time Pipeline & FR24 & TC24 & Section 5.5.4 \\ \hline
10. Urdu Localisation & FR25, FR26 & TC25 & Section 5.6.4 (partial) \\ \hline
11. Health Data Sync & FR27--FR29 & TC26, TC27 & Section 5.6.5 (not built) \\ \hline
12. AI Chat Coach & FR30--FR32 & TC28 & Section 5.5.8 (partial) \\ \hline
\end{tabularx}
\end{table}

\section{Evaluation Against Project Objectives}
\begin{table}[H]
\centering
\caption{Objectives and Results}
\label{tab:objectives}
\small
\begin{tabularx}{\textwidth}{|>{\raggedright\arraybackslash}p{5.2cm}|>{\raggedright\arraybackslash}p{2.6cm}|>{\raggedright\arraybackslash}X|}
\hline
\textbf{Objective (Section 1.6)} & \textbf{Result} & \textbf{Evidence} \\ \hline
Recognise 80+ foods with at least 85\% top-1 accuracy & Partly met & 217 classes (more than 80), but 58.5\% top-1 (Table~\ref{tab:modeleval}) \\ \hline
Detect several items in one image & Not met & The built model is a classifier (TC14) \\ \hline
Return nutrition data within 3 seconds & Partly met & Met on the YOLOv8 path, not on the Claude path (TC10) \\ \hline
Estimate portions from one image with no reference object & Partly met & Relative bucket only; no calibrated grams \\ \hline
Convert calories into activity minutes (MET) & Partly met & Used for logged activity; the meal-to-minutes display is not built \\ \hline
English and Urdu interface with right-to-left layout & Not met & Only Roman-Urdu greetings (TC25) \\ \hline
Complete app with camera, results, history and dashboard & Met on the simulator & TC09, TC11, TC12, TC23; camera pending a phone test \\ \hline
AI coach with personal advice using RAG & Not met & The chat needs a key and has no server-side retrieval (TC28) \\ \hline
Two-way sync with Google Fit and Apple HealthKit & Not met & Not built (TC26, TC27) \\ \hline
Quantised model with fast inference & Met & Static INT8: 3.6 times smaller, as fast, $-3.7$ points top-1 (Table~\ref{tab:quant}) \\ \hline
\end{tabularx}
\end{table}

\section{Limitations}
\begin{itemize}[noitemsep]
  \item The system names one dish per photo. The live mode sends one frame at a time and is not continuous
        video analysis.
  \item The model knows only the 217 classes it was trained on, and its top-1 accuracy is 58.5\%. Dishes outside this
        list, and similar-looking dishes, may be named wrongly, so the user must then choose the right item.
  \item Portion size is a relative estimate from one image. Camera angle, light, and very deep or heaped dishes can
        reduce its accuracy.
  \item Mixed dishes and hidden ingredients (oil, ghee, sugar) are hard to see, so calories can differ from the
        real value.
  \item Recognition needs an internet connection; offline logging is not built.
  \item Burned calories are entered by hand and estimated from a MET value; they are not measured.
  \item The quantised model is not yet running on the phone, and its accuracy was measured on a rebuilt split that may
        overlap the training data.
  \item The API does not yet verify Firebase tokens, and the chat coach key is held in the app.
  \item The system tracks calories and macronutrients only. Vitamins and minerals are not included.
  \item The coach depends on an external service and cannot replace a doctor or dietitian.
\end{itemize}

\section{Conclusion}
Pakalorie was created to solve a clear problem: most calorie tracking apps do not understand Pakistani food, and
typing every meal makes people stop using them. The project responds with a mobile app that turns one photo into
calories and macronutrients in a few seconds.

Nine of the twelve modules are built. A YOLOv8 classifier (or a Claude vision option) names the dish, MiDaS gives a
relative portion bucket, a Pakistani and USDA nutrition database supplies the values, and a calorie engine turns these
into results grounded in real database rows. Meal history, trends, a macro split and a calorie compensation card with
an activity log help users follow their goals over time. The recognition model was quantised to a file 3.6 times
smaller with a small loss of accuracy, and a live mode shows predictions from the camera through the server.

The work also shows its limits. The recognition accuracy of 58.5\% is well below the 85\% target, so the accuracy
objective is only partly met. Urdu localisation, health sync and a server-side AI coach are not complete, and several
tests are still to be run on a real phone. These are reported openly in the tables above, so that the next milestone
can start from a clear picture of what works.

\section{Future Work}
\begin{itemize}[noitemsep]
  \item Add more Pakistani and regional dishes and more images per class, and train a larger model, to raise accuracy
        on mixed and similar-looking dishes.
  \item Rebuild the test split so the accuracy of the quantised model can be reported on unseen images.
  \item Run the static INT8 model on the phone for live recognition without the server.
  \item Run the recognition and depth models in parallel on the server, as in the original design.
  \item Replace the manual activity entry with automatic data from Google Fit and Apple Health (Health Data Sync).
  \item Complete Urdu localisation with Urdu script and right-to-left layout.
  \item Move the AI chat coach to the server with retrieval, and verify Firebase tokens on every request.
  \item Separate foods on a mixed plate more finely (multi-dish plate decomposition).
  \item Improve portion estimation, for example with better depth models or the phone's depth sensor where it exists.
  \item Track micronutrients such as vitamins and minerals, and add barcode scanning for packaged products.
  \item Run the usability study with the SUS questionnaire and report the results.
\end{itemize}
